\section{Methodology}

\subsection{State Definition}

The implemented nominal state is the 2D NavState matrix
\begin{equation}
X =
\begin{bmatrix}
R(\theta) & p & v \\
0_{1\times 2} & 1 & 0 \\
0_{1\times 2} & 0 & 1
\end{bmatrix},
\label{eq:navstate}
\end{equation}
where \(R(\theta)\in \SO(2)\), \(p=[x\ y]^T\in\R^2\), and \(v=[\dot{x}\ \dot{y}]^T\in\R^2\). This state has five degrees of freedom:
\begin{equation}
    x_{\mathrm{dof}} =
    \begin{bmatrix}
    \theta & x & y & \dot{x} & \dot{y}
    \end{bmatrix}^{T}.
\end{equation}
The covariance \(P\in\R^{5\times 5}\) is not stored over the matrix entries of \(X\). It is stored over local tangent coordinates ordered as
\begin{equation}
    \delta x =
    \begin{bmatrix}
    \delta \theta &
    \delta p_x &
    \delta p_y &
    \delta v_x &
    \delta v_y
    \end{bmatrix}^{T}.
\label{eq:error_state}
\end{equation}

The helper \texttt{make\_navstate.m} constructs Eq.~(\ref{eq:navstate}); \texttt{navstate\_components.m} extracts \(\theta\), \(p\), and \(v\). The implementation uses a 2D analog of the GTSAM NavState idea, where pose and velocity are treated as one navigation state on a smooth group-like space \citep{gtsamNavigation,gtsamNavState,forster2015imu}.

\subsection{NavState Exponential Map}

Prediction and correction increments are mapped back to the nominal state by \texttt{navstate\_exp.m}. For a local increment
\begin{equation}
    \delta =
    \begin{bmatrix}
    \phi & \rho_x & \rho_y & \nu_x & \nu_y
    \end{bmatrix}^{T},
\end{equation}
define \(\rho=[\rho_x\ \rho_y]^T\), \(\nu=[\nu_x\ \nu_y]^T\), and
\begin{equation}
    A(\phi)=1-\frac{\phi^2}{6}, \qquad
    B(\phi)=\frac{\phi}{2}-\frac{\phi^3}{24}.
\end{equation}
The code uses the third-order Taylor left-Jacobian approximation
\begin{equation}
    J(\phi)=
    \begin{bmatrix}
    A(\phi) & -B(\phi) \\
    B(\phi) & A(\phi)
    \end{bmatrix}.
\end{equation}
The returned increment is
\begin{equation}
\Exp(\delta)=
\begin{bmatrix}
R(\phi) & J(\phi)\rho & J(\phi)\nu \\
0_{1\times 2} & 1 & 0 \\
0_{1\times 2} & 0 & 1
\end{bmatrix}.
\label{eq:nav_exp}
\end{equation}
The nominal update is group composition:
\begin{equation}
    X^{+}=X\,\Exp(\delta).
\end{equation}
Consequently, heading is updated on \(\SO(2)\), and position and velocity increments are rotated by the current nominal frame through matrix multiplication.

\subsection{Prediction Model}

The prediction input generated by \path{simulate_prediction_inputs.m} is body-frame acceleration and yaw rate:
\begin{equation}
    u_k =
    \begin{bmatrix}
    a_{b,k}^{T} & \omega_k
    \end{bmatrix}^{T}.
\end{equation}
The simulator computes true body-frame acceleration from the saved truth trajectory, then adds a scenario-constant bias and white noise. The filter does not estimate these biases. The configured generation noise is \(\sigma_a=0.08~\mathrm{m/s^2}\) and \(\sigma_\omega=0.012~\mathrm{rad/s}\), with per-scenario bias standard deviations \(0.018~\mathrm{m/s^2}\) and \(0.004~\mathrm{rad/s}\). The filter prediction covariance uses the larger values \(\sigma_a=0.11~\mathrm{m/s^2}\) and \(\sigma_\omega=0.018~\mathrm{rad/s}\).

At timestep \(k\), \texttt{propagate\_navstate.m} extracts \((\theta_k,p_k,v_k)\), forms \(R_k=R(\theta_k)\), and computes
\begin{align}
    \phi_k &= \omega_k \Delta t,\\
    \rho_k &= R_k^T(v_k\Delta t)+\frac{1}{2}a_{b,k}\Delta t^2,\\
    \nu_k &= a_{b,k}\Delta t.
\end{align}
The nominal propagation is
\begin{equation}
    X_{k+1|k}=X_{k|k}\Exp
    \left(
    \begin{bmatrix}
    \phi_k & \rho_k^T & \nu_k^T
    \end{bmatrix}^{T}
    \right).
\label{eq:prediction}
\end{equation}
For small \(\phi_k\), Eq.~(\ref{eq:prediction}) reduces to the familiar planar kinematic update \(p_{k+1}\approx p_k+v_k\Delta t+\frac{1}{2}R_k a_{b,k}\Delta t^2\), \(v_{k+1}\approx v_k+R_k a_{b,k}\Delta t\), and \(\theta_{k+1}\approx \theta_k+\omega_k\Delta t\).

The covariance propagation follows the EKF form \citep{simon2006optimal,brownHwang2012random}
\begin{equation}
    P_{k+1|k}=F_k P_{k|k}F_k^T+G_k Q_u G_k^T+Q_{\mathrm{floor}},
\label{eq:cov_prediction}
\end{equation}
where
\begin{equation}
    Q_u=\diag(\sigma_a^2,\sigma_a^2,\sigma_\omega^2)
\end{equation}
and
\begin{equation}
    Q_{\mathrm{floor}}=\diag(10^{-7},10^{-6},10^{-6},10^{-5},10^{-5}).
\end{equation}
Let
\begin{equation}
    S_2=
    \begin{bmatrix}
    0 & -1\\
    1 & 0
    \end{bmatrix},
    \qquad
    c_k=R_k S_2 a_{b,k}.
\end{equation}
The implemented linearization is
\begin{equation}
F_k =
\begin{bmatrix}
1 & 0_{1\times 2} & 0_{1\times 2}\\
\frac{1}{2}c_k\Delta t^2 & I_2 & \Delta t I_2\\
c_k\Delta t & 0_{2\times 2} & I_2
\end{bmatrix},
\label{eq:F}
\end{equation}
with block rows ordered as \(\theta,p,v\), and
\begin{equation}
G_k =
\begin{bmatrix}
0_{1\times 2} & \Delta t\\
\frac{1}{2}R_k\Delta t^2 & 0_{2\times 1}\\
R_k\Delta t & 0_{2\times 1}
\end{bmatrix}.
\label{eq:G}
\end{equation}

\subsection{Beacon Simulation Model}

For beacon \(i\), the true range and simulator-side hidden attenuation are
\begin{align}
    d_{i,k} &= \norm{p_k-b_i},\\
    A_{i,k} &= \sum_{j\in\mathcal{I}_{i,k}} t_j m_j,
\end{align}
where \(b_i\) is the known beacon position, \(\mathcal{I}_{i,k}\) is the set of wall segments intersected by the robot-to-beacon line segment, \(t_j\) is wall thickness, and \(m_j\) is material factor. This attenuation is computed by \texttt{wall\_attenuation.m} and is privileged simulator knowledge.

The simulator rejects a beacon measurement if any of the following conditions holds:
\begin{align}
    d_{i,k} &> d_{\max,i},&
    A_{i,k} &> A_{\max,i},&
    \mathrm{RSSI}_{i,k} &< \mathrm{RSSI}_{\min},
\end{align}
or if a random dropout draw succeeds. Before the RSSI threshold test, the implemented signal model is
\begin{equation}
    \mathrm{RSSI}_{i,k}
    =
    r_{0,i}
    -\alpha_d d_{i,k}
    -\alpha_w A_{i,k}
    + n_{s,k},
\label{eq:rssi}
\end{equation}
where \(\alpha_d=3.0\), \(\alpha_w=34.0\), \(n_{s,k}\sim\mathcal{N}(0,\sigma_s^2)\), and \(\sigma_s=2.2\). The repository stores a \texttt{path\_loss\_exponent} field for each beacon, but the implemented RSSI model is the linear distance-and-wall model in Eq.~(\ref{eq:rssi}), not a logarithmic path-loss model. The model is therefore a controlled synthetic degradation model rather than a full radio propagation model \citep{rappaport2002wireless}.

If the measurement remains available, the simulator uses
\begin{equation}
    \sigma_{r,i,k}
    =
    \sigma_{0}
    + k_d d_{i,k}
    + k_s\max(0,r_{0,i}-\mathrm{RSSI}_{i,k})
\end{equation}
with \(\sigma_0=0.12~\mathrm{m}\), \(k_d=0.025\), and \(k_s=0.006\). The measured range is
\begin{equation}
    z_{i,k}=d_{i,k}+n_{r,k},\qquad
    n_{r,k}\sim\mathcal{N}(0,\sigma_{r,i,k}^{2}).
\end{equation}

\subsection{Beacon EKF Update}

The filter-side beacon measurement function is
\begin{equation}
    h_i(X)=\norm{p-b_i}.
\end{equation}
The implementation clamps the predicted range used in the residual and Jacobian to the configured minimum \(d_{\min}=0.20~\mathrm{m}\):
\begin{equation}
    \hat{d}_{i,k}=\max\left(\norm{\hat{p}_k-b_i},d_{\min}\right).
\end{equation}
The scalar residual is
\begin{equation}
    r_{i,k}=z_{i,k}-\hat{d}_{i,k}.
\end{equation}
For the tangent ordering in Eq.~(\ref{eq:error_state}), the implemented beacon Jacobian is
\begin{equation}
H_{i,k}=
\begin{bmatrix}
0 &
\dfrac{\hat{p}_{x,k}-b_{x,i}}{\hat{d}_{i,k}} &
\dfrac{\hat{p}_{y,k}-b_{y,i}}{\hat{d}_{i,k}} &
0 & 0
\end{bmatrix}.
\label{eq:beacon_H}
\end{equation}

The fixed-covariance beacon case uses
\begin{equation}
    R_{i,k}=0.55^2.
\end{equation}
The adaptive beacon case uses only robot-side values, measured range and RSSI:
\begin{equation}
    \hat{\sigma}_{i,k}
    =
    \sigma_0
    + k_d\max(z_{i,k},0)
    + k_s\max(0,r_{0,i}-\mathrm{RSSI}_{i,k}),
\label{eq:adaptive_sigma}
\end{equation}
followed by the lower bound \(\hat{\sigma}_{i,k}\geq 0.05~\mathrm{m}\), and \(R_{i,k}=\hat{\sigma}_{i,k}^2\). The filter does not use true attenuation, true range before noise, wall counts, or wall thickness values during the beacon update.

\subsection{LiDAR Simulation and Pose Update}

For each scan, \texttt{raycast\_map.m} casts beams at angles
\begin{equation}
    \alpha_{\ell,k}=\theta_k+\beta_{\ell},
\end{equation}
where \(\beta_{\ell}\) spans the configured \(270^\circ\) field of view. The range for beam \(\ell\) is the closest wall intersection distance within 12 m, or the maximum range if no intersection exists. Valid returns are corrupted by independent Gaussian range noise with \(\sigma=0.035~\mathrm{m}\); invalid returns are stored at max range.

The scan quality score is computed from valid return fraction \(f_v\), capped normalized range spread \(s_r\), and near-wall fraction \(f_n\):
\begin{align}
    q_k=\operatorname{clamp}\big(
    &0.10+0.50f_v+0.25s_r+0.15f_n,\nonumber\\
    &0.05,\ 0.98
    \big).
\end{align}
For the open-hallway scenario, the code subtracts 0.18 from this score when the range spread is less than 0.16, modeling reduced geometric diversity in hallway-like scans.

At the LiDAR pose-update rate, the simulator forms
\begin{equation}
    z_{\ell,k} =
    \begin{bmatrix}
    x_k & y_k & \theta_k
    \end{bmatrix}^{T}
    + n_{\ell,k},
\end{equation}
when the quality and valid-return tests pass and no random dropout occurs. The covariance is
\begin{equation}
R_{\ell,k}=\diag
\left(
\sigma_{xy,k}^{2},\sigma_{xy,k}^{2},\sigma_{\theta,k}^{2}
\right),
\end{equation}
with
\begin{align}
    \sigma_{xy,k} &= 0.08+0.55(1-q_k)^2,\\
    \sigma_{\theta,k} &= 1.2^\circ+7.0^\circ(1-q_k)^2.
\end{align}
The EKF measurement model is
\begin{equation}
    h_{\ell}(X)=
    \begin{bmatrix}
    p_x & p_y & \theta
    \end{bmatrix}^{T},
\end{equation}
with residual
\begin{equation}
    r_{\ell,k}=
    \begin{bmatrix}
    z_x-\hat{p}_x\\
    z_y-\hat{p}_y\\
    \wrap(z_\theta-\hat{\theta})
    \end{bmatrix}
\end{equation}
and Jacobian
\begin{equation}
H_{\ell}=
\begin{bmatrix}
0 & 1 & 0 & 0 & 0\\
0 & 0 & 1 & 0 & 0\\
1 & 0 & 0 & 0 & 0
\end{bmatrix}.
\label{eq:lidar_H}
\end{equation}
This synthetic pose-measurement interface is a deliberate simplification consistent with probabilistic range-sensor modeling practice \citep{thrun2005probabilistic,dellaertHutchinson2023robotics}; no scan matching is implemented.

\subsection{EKF Correction}

Both beacon and LiDAR updates call \texttt{ekf\_tangent\_update.m}. For residual \(r_k\), Jacobian \(H_k\), and measurement covariance \(R_k\), the innovation covariance and Kalman gain are
\begin{align}
    S_k &= H_kP_{k|k-1}H_k^T+R_k,\\
    K_k &= P_{k|k-1}H_k^T S_k^{-1}.
\end{align}
The tangent correction is
\begin{equation}
    \delta x_k=K_k r_k.
\end{equation}
If \(\norm{\delta x_k}>4.0\), the code scales the correction to norm 4.0 before applying it. Let \(\delta p_k\) and \(\delta v_k\) denote the world-frame translational components of \(\delta x_k\). Before composition, the code converts those components into the local nominal frame:
\begin{equation}
    \delta_{\mathrm{group},k}=
    \begin{bmatrix}
    \delta \theta_k\\
    \hat{R}_k^T\delta p_k\\
    \hat{R}_k^T\delta v_k
    \end{bmatrix}.
\end{equation}
The nominal state update is
\begin{equation}
    X_{k|k}=X_{k|k-1}\Exp(\delta_{\mathrm{group},k}).
\end{equation}
The covariance uses the Joseph form
\begin{equation}
P_{k|k}
=
(I-K_kH_k)P_{k|k-1}(I-K_kH_k)^T
+K_kR_kK_k^T,
\end{equation}
followed by explicit symmetrization. The normalized innovation squared stored in the innovation history is
\begin{equation}
    \mathrm{NIS}_k=r_k^T S_k^{-1}r_k.
\end{equation}
